\section{Flat square-zero obstructions at every prime}
\label{sec:dual-numbers-all-primes}

A flat square-zero thickening can destroy generation in every residue
characteristic.  The obstruction is a boundary on periodic cohomology;
its square has primary components with additive lengths modulo two.

\begin{theorem}
\label{thm:dual-numbers-all-primes}
For every prime $p$, put
\[
 O_p=\mathbb Z_{(p)}\text{ or }\mathbb Z_p,\qquad
 A_p^{\mathrm{dual}}=O_p[\epsilon]/(\epsilon^2).
\]
There is an explicit $A_p^{\mathrm{dual}}[S_p]$-lattice $M_p$ of rank $p!$
whose class is not an integer combination of Young permutation classes.
Moreover
\[
 [M_p]-[A_p^{\mathrm{dual}}[S_p]]
\]
reduces to zero over $O_p$ and is detected by an additive
$\mathbb Z/2$-valued invariant.  In particular the polynomial map generated
by the trivial symmetric-group classes is not surjective over
$A_p^{\mathrm{dual}}$.
\end{theorem}

We give all details, including the effect of changing a generator of the
cyclic subgroup.  Proposition~\ref{prop:truncation-torsion} below
proves that the displayed class has exact order two for every prime;
Corollary~\ref{cor:all-prime-two-torsion} gives an infinite split family.

\subsection{Primary-length invariants for the cyclic group of order \texorpdfstring{$p$}{p}}

Fix $H=C_p=\langle s\rangle$, write $O=O_p$ and $A=A_p^{\mathrm{dual}}$,
and let $N$ be an $O[H]$-lattice with action matrix $S$.  Set
\[
 \mathcal N(S)=I+S+\cdots+S^{p-1}.
\]
Use the two-periodic complex with even differential $S-I$ and odd
differential $\mathcal N(S)$.  Its cohomology is
\[
 H^0(N)=\frac{\ker(S-I)}{\mathcal N(S)N},\qquad
 H^1(N)=\frac{\ker\mathcal N(S)}{(S-I)N}.
\]
Both groups are finite dimensional over $\mathbb F_p$.  For the even
group, this follows from $\mathcal N(S)v=pv$ on invariants.  For the odd
group use
\begin{equation}
\label{eq:dual-prime-D-identity}
 (S-I)D_p(S)=\mathcal N(S)-pI,
 \qquad
 D_p(t)=\sum_{j=0}^{p-2}(p-1-j)t^j.
\end{equation}
If $\mathcal N(S)v=0$, this identity expresses $pv$ as a boundary.

For an $A[H]$-lattice $L$, the canonical extension
\[
 0\longrightarrow L_0\xrightarrow{\epsilon}L\longrightarrow L_0
 \longrightarrow0,\qquad L_0=L/\epsilon L,
\]
gives a degree-one boundary $\delta_L$ on this periodic cohomology and an
even endomorphism $u_L=\delta_L^2$.  For every monic irreducible
$f\in\mathbb F_p[t]$ with $f\ne t$, define
\[
 \nu_f(L)=
 \ell_{\mathbb F_p[t]}
 \bigl(H^0(L_0)_{f\text{-primary}},\ t\text{ acting as }u_L\bigr)
 \pmod2.
\]
We prove additivity directly, including the sign issue when $p$ is odd.

For an admissible exact sequence $0\to L'\to L\to L''\to0$, choose an
$A$-linear splitting.  In each cochain degree write its differential as
\[
 d_L^j=\begin{pmatrix}d_{L'}^j&B^j\\0&d_{L''}^j\end{pmatrix},
\]
and expand every block into its constant and $\epsilon$ coefficients.
The coefficient of $\epsilon$ in the upper-right block of
$d_L^{j+1}d_L^j=0$ is
\begin{align*}
 d_{L',1}^{j+1}B_0^j+B_0^{j+1}d_{L'',1}^j
 ={}&-d_{L',0}^{j+1}B_1^j-B_1^{j+1}d_{L'',0}^j.
\end{align*}
On a degree-$j$ cocycle for $L''_0$, the last term is zero and the first
term on the right is a boundary.  The left side represents
\[
 \delta_{L'}^{j+1}\partial_j+
 \partial_{j+1}\delta_{L''}^j.
\]
Thus the extension boundaries anticommute with the long-exact-sequence
boundaries, and all maps commute with $u=\delta^2$.

Taking an $f$-primary component is exact, as is seen by a common finite
Bezout projector.  Since $f\ne t$, $u$ is invertible there, and $\delta$
identifies the two parities by $\mathbb F_p[t]$-linear isomorphisms.  The
two boundary images in the resulting exact six-term cycle have equal
composition length.  The minus sign in their intertwining identity is a
unit and does not affect that equality.  If $a,b,c$ are the three
degree-zero lengths and $r_0,r_1$ are the two boundary-image lengths, then
\[
 a-b+c=r_0+r_1=2r_0.
\]
This proves that $\nu_f$ is a homomorphism on
$K_0(\operatorname{Lat}_A(H))$.  It vanishes on every lattice obtained by
scalar extension from $O$, because its canonical $\epsilon$ extension then
splits as an $O[H]$-extension and $\delta=0$.

\begin{proposition}\label{prop:dual-bockstein-invariants}
For every monic irreducible $f\in\mathbb F_p[t]$ other than $t$, the
primary-length formula above defines a homomorphism
\[
 \nu_f:K_0(\operatorname{Lat}_{O[\epsilon]/(\epsilon^2)}(C_p))
       \longrightarrow\mathbb Z/2.
\]
It vanishes on scalar extensions from $O$.  For any fixed class only
finitely many $\nu_f$ are nonzero.
\end{proposition}
The final assertion follows from finite dimension of the cohomology and
then from finite support of a finite integer combination of lattice classes.


\subsection{An explicit rank-\texorpdfstring{$p$}{p} cyclic lattice}

Put $r=p-1$.  Let $C$ be the integral companion matrix of
$\Phi_p(t)=1+t+\cdots+t^{p-1}$, with basis $e_1,\ldots,e_r$ and action
\[
 Ce_i=e_{i+1}\ (i<r),\qquad
 Ce_r=-\sum_{i=1}^r e_i.
\]
Thus $C^p=I$ and $\mathcal N(C)=0$.  Set
\[
 b=e_1,\qquad a=e_r^{\mathsf T},\qquad
 T_p=\begin{pmatrix}1&\epsilon a\\\epsilon b&C\end{pmatrix}.
\]
The diagonal blocks of $T_p^p$ are $1$ and $C^p$, since every off-diagonal
round trip contains $\epsilon^2$.  Its off-diagonal blocks are respectively
\[
 \epsilon a\mathcal N(C),\qquad
 \epsilon\mathcal N(C)b,
\]
and therefore vanish.  Hence $T_p^p=I$, giving an $A[H]$-lattice
$L_p=A^p$.

The reduced representation is $O\oplus I$, where the companion module $I$
is isomorphic to the augmentation ideal of $O[H]$: identify $e_i$ with
$s^{i-1}(s-1)$ for $1\le i\le p-1$.  The map $C-I$ is injective over $O$
and has cokernel $\mathbb F_p$, with the common class of all $e_i$ as a
generator.  Therefore
\[
 H^0((L_p)_0)=\mathbb F_p\quad\text{from the trivial summand},
 \qquad
 H^1((L_p)_0)=\mathbb F_p\quad\text{from }I.
\]

The even extension boundary sends the trivial generator to the class of
$b=e_1$, hence is the identity in these coordinates.  For the odd boundary,
the upper-right $\epsilon$ coefficient of $\mathcal N(T_p)$ is
$aD_p(C)$.  Since $C^jb=e_{j+1}$ for $0\le j\le p-2$,
\[
 aD_p(C)b=1.
\]
Thus the odd boundary is also the identity.  We conclude that
\begin{equation}
\label{eq:dual-prime-u-one}
 u_{L_p}=1\quad\text{on both cohomology spaces},\qquad
 \nu_{t-1}(L_p)=1.
\end{equation}

\subsection{Changing the generator: exact chain maps and the scaling factor}

For $1\le q\le p-1$, let $L^{(q)}$ denote the twist in which $s$ acts as
$T^q$ rather than $T$.  Put
\[
 S_q(T)=I+T+\cdots+T^{q-1}.
\]
There is a cochain map from the periodic complex for $T$ to the periodic
complex for $T^q$, in the following direction and degrees:
\begin{equation}
\label{eq:dual-prime-comparison}
 F_q^{2j}=q^j I,\qquad
 F_q^{2j+1}=q^j S_q(T),\qquad j\in\mathbb Z.
\end{equation}
Negative powers of $q$ are allowed because $q\in O^\times$.  The two chain
identities are
\[
 T^q-I=S_q(T)(T-I),\qquad
 \mathcal N(T^q)S_q(T)=\mathcal N(T)S_q(T)=q\mathcal N(T).
\]
Here $\mathcal N(T^q)=\mathcal N(T)$ because multiplication by $q$ permutes
the residues modulo $p$.

The same maps apply before and after reduction modulo $\epsilon$ and
commute with the canonical extension.  On degree-zero cohomology,
$F_q^0=I$ identifies the two groups: $T$ and $T^q$ have the same invariants
and the same norm image.  On degree two, however, the map is $F_q^2=qI$.
Naturality of the extension boundaries gives
\[
 \delta_{L^{(q)}}^0F_q^0=F_q^1\delta_L^0,
 \qquad
 \delta_{L^{(q)}}^1F_q^1=F_q^2\delta_L^1.
\]
Composing and using the canonical period-two identification yields
\begin{equation}
\label{eq:dual-prime-scaling}
 u_{L^{(q)}}=q\,u_L
 \quad\text{on the commonly identified }H^0.
\end{equation}
This is why it would be incorrect to treat the periodicity identification
as invariant under the generator change without its scalar factor.

For $L=L_p$, equations \eqref{eq:dual-prime-u-one} and
\eqref{eq:dual-prime-scaling} give
\begin{equation}
\label{eq:dual-prime-twist-detector}
 \nu_{t-1}(L_p^{(q)})=
 \begin{cases}1,&q=1\text{ in }\mathbb F_p,\\0,&q\ne1.
 \end{cases}
\end{equation}

\subsection{Induction to the symmetric group and the Mackey multiplicities}

Embed $H$ in $S_p$ as the subgroup generated by a $p$-cycle and define
\[
 M_p=\operatorname{Ind}_H^{S_p}L_p.
\]
Its $A$-rank is $[S_p:H]p=p!$.  To distinguish its class, apply
$\nu_{t-1}$ after restriction back to $H$.

For a double coset in $H\backslash S_p/H$, the intersection
$H\cap xHx^{-1}$ is either $H$ or the trivial group, since $H$ has prime
order.  The first case occurs exactly when $x\in N_{S_p}(H)$.  For every
other double coset, the Mackey summand is induced from the trivial
subgroup, hence is free over $A[H]$.  Its reduced periodic cohomology is
zero, so it contributes zero to every $\nu_f$.

Inside the normalizer, the double cosets are exactly its cosets modulo
$H$.  The centralizer of a $p$-cycle in $S_p$ is $H$, and every automorphism
$s\mapsto s^q$, $q\in\mathbb F_p^\times$, is realized by a permutation of
the $p$ cyclically labeled letters.  Thus
\[
 N_{S_p}(H)/H\cong\mathbb F_p^\times.
\]
The remaining Mackey summands are precisely one copy of each
$L_p^{(q)}$, with $q$ running through the nonzero residues.  Changing a
Mackey convention replaces $q$ by $q^{-1}$ and does not change either the
multiplicity statement or the unique occurrence of $q=1$.  Consequently
\[
 \nu_{t-1}\!\left(\operatorname{Res}^{S_p}_H M_p\right)
 =\sum_{q\in\mathbb F_p^\times}\nu_{t-1}(L_p^{(q)})=1.
\]

Every Young permutation lattice over $A$ is obtained by scalar extension
from its integral $O$-model.  Its restriction to $H$ is also such a scalar
extension, so its invariant is zero.  Therefore $[M_p]$ is not in their
integer span.

Finally, $[(L_p)_0]=[O]+[I]=[O[H]]$ by the augmentation exact sequence.
Induction gives $[(M_p)_0]=[O[S_p]]$.  After inverting $p$, the matrix
$C-I$ is invertible and a change of basis $I+\epsilon X$, with
$X_{12}=a(C-I)^{-1}$ and $X_{21}=(I-C)^{-1}b$, identifies $L_p$ with its
constant reduction.  That reduction over $\operatorname{Frac}(O)$ is the regular
cyclic representation.  Thus the same equality holds after inverting
$p$, even with $\epsilon$ retained.  The class
$[M_p]-[A[S_p]]$ is thus in the asserted reduction kernel and is nonzero by
the displayed invariant.  This completes the proof of
Theorem~\ref{thm:dual-numbers-all-primes}.
